The Model Context Protocol (MCP) represents a significant advancement in addressing one of the most persistent challenges in AI agent architecture: the effective management and sharing of context across interactions and agent boundaries. This section examines the origins, architecture, implementation approaches, and specific advantages of MCP, with particular focus on how it addresses the disconnected models problem identified in previous research.\\
\subsection{3.1 Origins and Development}
\subsubsection{Historical Context and Development Timeline}
The Model Context Protocol emerged in response to a growing recognition of the limitations of existing approaches to context management in AI systems. As large language models became increasingly capable, their integration into agent architectures revealed a fundamental disconnect between the sophisticated reasoning capabilities of these models and their ability to maintain coherent context across interactions.\\ \newline
The development of MCP can be traced through several key milestones:\\ \newline
**Late 2023 - Early 2024**: Initial research and conceptualization at Anthropic, driven by challenges encountered in developing Claude and related agent systems. Internal prototypes demonstrated the potential benefits of standardized context management.\\ \newline
**Mid 2024**: Formal introduction of MCP as an open protocol, with Anthropic releasing initial specifications and reference implementations. The protocol was explicitly positioned as an open standard rather than a proprietary solution, reflecting recognition of the need for industry-wide collaboration.\\ \newline
**Late 2024**: Public release of MCP 1.0, including comprehensive documentation, SDKs for multiple programming languages, and sample implementations for common use cases. This release established the core architecture and primitives that define the protocol.\\ \newline
**Early 2025**: Rapid adoption across the AI ecosystem, with multiple organizations implementing MCP-compatible systems and contributing extensions for specialized domains. Formation of an informal working group to guide protocol evolution and standardization efforts.\\ \newline
**Present (Mid 2025)**: Ongoing refinement and extension of the protocol, with particular focus on remote/cloud integration, enhanced security models, and support for additional modalities beyond text.\\ \newline
This development trajectory reflects a deliberate effort to address a critical industry need through open collaboration rather than proprietary solutions. By establishing MCP as an open standard with clear specifications and reference implementations, Anthropic and subsequent contributors have created a foundation for interoperability that transcends individual platforms or vendors.\\
\subsubsection{Design Principles and Objectives}
The development of MCP has been guided by several core design principles that shape its architecture and implementation:\\ \newline
1. **Interoperability**: MCP is designed to work across different AI models, platforms, and environments, enabling consistent context management regardless of the specific technologies involved. This principle is reflected in the protocol's language-agnostic design and use of standard data formats.\\ \newline
2. **Simplicity**: The protocol prioritizes simplicity and ease of implementation, focusing on a minimal set of core primitives that can be combined to support complex use cases. This approach reduces barriers to adoption and encourages consistent implementation across diverse systems.\\ \newline
3. **Extensibility**: While maintaining a simple core, MCP is designed to be extensible through well-defined mechanisms for adding new capabilities and adapting to specialized domains. This balance between stability and evolution ensures the protocol can address emerging needs without sacrificing compatibility.\\ \newline
4. **Security and privacy by design**: MCP incorporates security and privacy considerations as fundamental design elements rather than afterthoughts. This includes clear permission models, data minimization principles, and mechanisms for controlling information flow between components.\\ \newline
5. **Human-centered control**: The protocol is designed to maintain appropriate human oversight and control, particularly for sensitive operations or decisions. This principle reflects recognition of the importance of human judgment in AI systems.\\ \newline
These design principles support MCP's primary objectives:\\ \newline
- Enabling AI models to access external information and tools in a standardized way\\ \newline
- Facilitating context sharing across agent boundaries and interaction sessions\\ \newline
- Reducing implementation complexity for developers integrating AI with external systems\\ \newline
- Establishing a foundation for interoperability across the AI ecosystem\\ \newline
- Enhancing agent capabilities through consistent access to contextual information\\ \newline
By adhering to these principles and objectives, MCP provides a robust foundation for addressing the context management challenges that have limited agent effectiveness in previous systems.\\
\subsubsection{Relationship to Previous Context Management Approaches}
MCP builds upon and extends several previous approaches to context management in AI systems, incorporating lessons learned while addressing limitations:\\ \newline
**Prompt engineering techniques**, which embed contextual information directly within model inputs, provided a simple but limited approach to context management. While effective for short interactions, these techniques face inherent limitations due to context window constraints and lack structured mechanisms for context prioritization or persistence. MCP incorporates the flexibility of prompt-based approaches while providing more systematic management of what context is included and when.\\ \newline
**Retrieval-augmented generation (RAG)** systems demonstrated the value of dynamically incorporating relevant information from external sources into model contexts. However, most RAG implementations focus narrowly on document retrieval rather than broader context management, and lack standardized interfaces for diverse information sources. MCP generalizes the RAG concept to encompass a wider range of context types and sources, with standardized interfaces that simplify integration.\\ \newline
**Tool-using agent frameworks** like LangChain and AutoGPT established patterns for connecting models with external tools and APIs, enabling more capable agent behaviors. These frameworks typically implement custom, framework-specific approaches to tool integration and context management. MCP standardizes these integration patterns, providing consistent interfaces that work across frameworks and environments.\\ \newline
**Memory management systems** in agent architectures have explored various approaches to storing and retrieving interaction histories and other contextual information. These systems often employ specialized memory types (working memory, episodic memory, semantic memory) but typically lack standardized interfaces or interoperability. MCP provides a unified framework for implementing diverse memory types with consistent access patterns.\\ \newline
By synthesizing insights from these previous approaches while addressing their limitations, MCP represents a significant advancement in context management for AI systems. Its standardized, interoperable design overcomes many of the integration challenges that have limited the effectiveness of earlier approaches.\\
\subsection{3.2 Architecture and Components}
\subsubsection{Client-Server Architecture}
MCP implements a client-server architecture that cleanly separates AI models (clients) from the data sources and tools they access (servers). This architectural pattern provides several advantages:\\ \newline
1. **Clear separation of concerns**: The client-server division establishes distinct responsibilities, with clients focusing on AI reasoning and decision-making while servers handle data access and tool execution.\\ \newline
2. **Flexible deployment options**: Clients and servers can be deployed in various configurations—on the same machine, within the same network, or across remote environments—depending on security requirements and operational constraints.\\ \newline
3. **Scalability**: The architecture supports both one-to-many relationships (a single client accessing multiple servers) and many-to-one relationships (multiple clients sharing access to common servers), enabling flexible scaling patterns.\\ \newline
4. **Independent evolution**: Clients and servers can evolve independently as long as they maintain compatibility with the protocol, allowing for innovation on both sides without requiring lockstep development.\\ \newline
In a typical MCP implementation, the client component is integrated with an AI application (such as an agent framework or assistant interface) and manages connections to one or more MCP servers. Each server provides access to specific data sources or tools, exposing their capabilities through standardized MCP interfaces.\\ \newline
Communication between clients and servers follows a request-response pattern using JSON-RPC, a lightweight remote procedure call protocol that uses JSON for data encoding. This approach provides a language-agnostic interface that can be implemented across diverse platforms and environments.\\
\subsubsection{Standardized Primitives}
MCP defines a set of core primitives—standardized operations and data structures—that provide the building blocks for context management. These primitives are divided into server-side and client-side categories:\\ \newline
**Server-side primitives** define the capabilities that MCP servers expose to clients:\\ \newline
1. **Prompts**: Pre-defined instructions or templates that an AI can request. These may include situational guidance, formatting instructions, or specialized procedures for particular tasks. Prompts help establish consistent behavior patterns and can reduce the need to repeatedly include standard instructions in the model's context.\\ \newline
2. **Resources**: Structured data or documents that can be sent to the AI model. Resources may include knowledge base articles, database records, configuration information, or other contextual data that informs model responses. The resource primitive includes metadata that helps the client determine how and when to incorporate the resource into the model's context.\\ \newline
3. **Tools**: Executable functions that the AI can invoke to perform actions or retrieve information. Tools may include API calls, database queries, file operations, or other functional capabilities. Each tool is defined with a name, description, parameter schema, and return type, enabling the model to understand its purpose and usage.\\ \newline
**Client-side primitives** define capabilities that MCP clients provide to servers:\\ \newline
1. **Roots**: Entry points that give servers access to specific data realms on the client side. Roots establish permission boundaries, allowing clients to control which data each server can access. For example, a client might provide a file system root that grants access to a specific directory but not the entire file system.\\ \newline
2. **Sampling**: A mechanism that allows an MCP server to request the AI model to generate a completion. This powerful capability enables servers to leverage the model's reasoning abilities as part of fulfilling a request, essentially allowing the server to "ask back" when additional context or decisions are needed.\\ \newline
These primitives provide a comprehensive foundation for context management while maintaining a relatively simple interface. By combining these basic operations in different ways, MCP can support sophisticated context management patterns without requiring complex protocol extensions.\\
\subsubsection{Communication Mechanisms and Message Formats}
MCP communication is structured around a well-defined message format based on JSON-RPC 2.0, a lightweight remote procedure call protocol. This format provides a standardized way to encode requests, responses, and notifications between clients and servers.
The basic structure of an MCP request includes:

\begin{verbatim}

{
  "jsonrpc": "2.0",
  "id": "request-123",
  "method": "resource.get",
  "params": {}
    "resource_id": "document-456",
    "format": "markdown"
}
And a corresponding response:
{
  "jsonrpc": "2.0",
  "id": "request-123",
  "result": {}
    "content": "# Document Title
    This is the content of the requested document...",
    "metadata": {
      "author": "Jane Smith",
      "created_at": "2025-03-15T14:30:00Z"
    }
}

\end{verbatim}

This structured format ensures consistent interpretation across different implementations while remaining human-readable for debugging and development purposes.\\ \newline
MCP defines specific methods for each primitive type, such as:\\ \newline
- `prompt.list` and `prompt.get` for accessing prompts\\ \newline
- `resource.list` and `resource.get` for retrieving resources\\ \newline
- `tool.list`, `tool.describe`, and `tool.execute` for discovering and using tools\\ \newline
- `root.list` and `root.describe` for exploring available roots\\ \newline
- `sample.generate` for requesting model completions\\ \newline
Each method has defined parameter and return schemas that ensure consistent behavior across implementations. The protocol also includes mechanisms for error handling, allowing servers to return structured error information when requests cannot be fulfilled.\\ \newline
Beyond the basic request-response pattern, MCP supports more complex interaction patterns through features like:\\ \newline
1. Streaming responses: For large resources or long-running tool executions, servers can stream results incrementally rather than waiting for complete execution.\\ \newline
2. Notifications: One-way messages that don't require responses, useful for events or updates that don't need acknowledgment.\\ \newline
3. Batched requests: Multiple requests combined into a single message to reduce communication overhead.\\ \newline
These communication mechanisms provide the flexibility needed to support diverse context management scenarios while maintaining protocol simplicity and consistency.\\
\subsection{Implementation Approaches}
\subsubsection{Server-Side Implementation Patterns}
Implementing an MCP server involves exposing data sources or tools through the standardized MCP interface. Several common implementation patterns have emerged:\\ \newline
Adapter pattern: Many MCP servers function as adapters that translate between MCP's standardized interface and existing APIs or data sources. This pattern is particularly common for integrating with established systems like document repositories, databases, or SaaS platforms. The adapter handles authentication, request formatting, and response translation, presenting a consistent MCP interface regardless of the underlying system's native API.\\ \newline
Composite pattern: Some MCP servers aggregate multiple data sources or tools behind a unified interface. These composite servers may provide integrated search across multiple repositories, coordinated access to related tools, or domain-specific combinations of capabilities. This pattern simplifies client configuration by reducing the number of server connections needed while providing more sophisticated integration capabilities.\\ \newline
Proxy pattern: In environments with strict security requirements, MCP servers may act as proxies that enforce access controls, audit usage, or transform requests and responses for compliance purposes. These proxy servers sit between clients and the actual data sources or tools, providing an additional layer of control and monitoring.\\ \newline
Embedded pattern: For simple use cases or development scenarios, MCP servers may be embedded directly within applications rather than deployed as separate services. This pattern reduces deployment complexity but may limit scalability and reuse across multiple clients.\\ \newline
Regardless of the implementation pattern, effective MCP servers typically share several characteristics:\\ \newline
1. Clear capability boundaries: Well-defined scope of what the server provides, avoiding overly broad or ambiguous functionality.\\ \newline
2. Comprehensive metadata: Detailed descriptions of available prompts, resources, and tools to enable effective discovery and usage.\\ \newline
3. Robust error handling: Clear error messages and appropriate status codes when requests cannot be fulfilled.\\ \newline
4. Efficient resource management: Appropriate caching, connection pooling, and resource cleanup to ensure performance and reliability.\\ \newline
5. Security-conscious design: Careful attention to authentication, authorization, and data protection throughout the implementation.\\ \newline
The open-source ecosystem around MCP has produced reference implementations and libraries for common server patterns, reducing the implementation effort required for new integrations.\\
\subsubsection{Client-Side Integration Strategies}
On the client side, integrating MCP involves connecting AI models with MCP servers and managing the flow of contextual information. Several integration strategies have proven effective:\\ \newline
Framework integration: Many agent frameworks and AI development platforms now include built-in MCP client capabilities. These integrations typically provide high-level abstractions that simplify server discovery, connection management, and context incorporation. Developers using these frameworks can leverage MCP capabilities with minimal additional code, often through declarative configuration rather than explicit implementation.\\ \newline
SDK-based integration: For applications not using MCP-enabled frameworks, software development kits (SDKs) provide libraries and utilities for implementing MCP client functionality. These SDKs, available for multiple programming languages, handle the details of protocol implementation while providing idiomatic interfaces for each language ecosystem. This approach requires more implementation effort than framework integration but offers greater flexibility and control.\\ \newline
Custom implementation: Some specialized applications implement MCP client functionality directly, particularly when operating in constrained environments or when requiring optimizations not available in existing libraries. While this approach requires the most development effort, it allows for precise tailoring to specific requirements and may enable performance optimizations not possible with general-purpose implementations.\\ \newline
Regardless of the integration strategy, effective MCP clients typically implement several key functions:\\ \newline
1. Server discovery and connection management: Finding available MCP servers, establishing connections, and handling authentication and reconnection as needed.\\ \newline
2. Context selection and prioritization: Determining which contextual information is most relevant for the current interaction and managing context window constraints.\\ \newline
3. Tool invocation and result processing: Translating model outputs into tool calls and incorporating tool results back into the model's context.\\ \newline
4. Error handling and fallback strategies: Gracefully handling server unavailability, request failures, or other exceptional conditions.\\ \newline
5. Context persistence: Maintaining relevant context across multiple interactions or sessions, potentially using client-side storage for efficiency.\\ \newline
The choice of integration strategy depends on factors including the application's architecture, performance requirements, existing technology stack, and development resources. The flexibility of MCP allows for implementation across a wide range of scenarios, from simple prototypes to enterprise-scale deployments.\\
\subsubsection{Security and Authentication Considerations}
Security is a fundamental consideration in MCP implementation, particularly as the protocol often provides access to sensitive data or powerful capabilities. Several security dimensions require attention:\\ \newline
Authentication and authorization: MCP implementations must verify the identity of connecting clients and servers (authentication) and determine what operations they are permitted to perform (authorization). The protocol supports various authentication mechanisms, from simple API keys for development environments to more robust OAuth flows or mutual TLS for production systems. Authorization is typically implemented through capability-based security models, where access is granted to specific resources or tools rather than through broad permissions.\\ \newline
Data protection: Information transmitted through MCP may include sensitive data that requires protection both in transit and at rest. The protocol mandates TLS encryption for all communications, ensuring transport-level security. Additionally, implementations may employ end-to-end encryption for particularly sensitive data, ensuring that even the transport infrastructure cannot access unencrypted content.\\ \newline
Isolation and sandboxing: Tool execution in MCP servers presents potential security risks, particularly for operations that interact with underlying systems. Robust implementations employ sandboxing techniques to isolate tool execution environments, limiting the potential impact of malicious inputs or unexpected behaviors.\\ \newline
Audit and monitoring: Comprehensive logging and monitoring are essential for security management in MCP deployments. This includes recording access patterns, tracking resource usage, and monitoring for anomalous behaviors that might indicate security issues. The protocol includes standardized fields for request identification and correlation, facilitating effective audit trails across system boundaries.\\ \newline
Permission models: MCP implements fine-grained permission models through the root primitive and capability-based access controls. These models allow clients to precisely control what servers can access and what operations they can perform, following the principle of least privilege. For example, a file system root might grant access only to specific directories, while a tool might be restricted to particular parameter ranges or operation types.\\ \newline
These security considerations are not merely implementation details but fundamental aspects of the protocol design. By incorporating security and authentication as core protocol elements, MCP ensures that implementations can be both powerful and secure, addressing the critical requirements of enterprise and sensitive consumer applications.\\
\subsection{MCP as a Solution to the Disconnected Models Problem}
\subsubsection{Addressing Context Discontinuity Across Agent Actions}
The "disconnected models problem" identified by Krishnan (2025) and emphasized by Microsoft's Sam Schillace represents one of the most significant limitations in current agent architectures. As Schillace noted, "To be autonomous you have to carry context through a bunch of actions, but the models are very disconnected and don't have continuity the way we do." This discontinuity manifests as agents losing track of previous interactions, forgetting relevant information, or failing to maintain coherent reasoning across multiple steps.\\ \newline
MCP directly addresses this challenge through several mechanisms:\\ \newline
1. Standardized context storage and retrieval: MCP provides consistent patterns for storing contextual information externally and retrieving it when needed. This allows context to persist beyond the limitations of model context windows and across multiple interactions.\\ \newline
2. Context-aware tool execution: When agents invoke tools through MCP, the tool execution can incorporate relevant context from previous interactions. This ensures that actions build upon previous steps rather than occurring in isolation.\\ \newline
3. Shared context across agent boundaries: In multi-agent systems, MCP enables different agents to access the same contextual information, ensuring consistent understanding and coordinated action even when tasks span multiple specialized agents.\\ \newline
4. Contextual prioritization: MCP's resource primitive includes metadata that helps determine the relevance and importance of different contextual elements, enabling more intelligent management of limited context space.\\ \newline
These capabilities directly address the discontinuity problem, enabling agents to maintain coherent context across extended interaction sequences and complex tasks. By providing external context management that complements the model's internal processing, MCP overcomes the inherent limitations of context windows while preserving the flexibility and reasoning capabilities of large language models.\\
\subsubsection{Enabling Persistent Memory and Context Awareness}
Beyond addressing immediate context discontinuity, MCP enables more sophisticated forms of persistent memory and context awareness that enhance agent capabilities:\\ \newline
Episodic memory stores records of specific interactions or experiences, allowing agents to recall and reference previous encounters. MCP supports episodic memory through resource primitives that can capture interaction histories in structured formats, making them available for future reference. This capability is particularly valuable for maintaining relationship context in assistant applications or tracking progress in extended problem-solving scenarios.\\ \newline
Semantic memory organizes conceptual knowledge and learned information, enabling agents to build and refine their understanding over time. MCP facilitates semantic memory through resources that represent structured knowledge, potentially including knowledge graphs, concept maps, or other semantic representations. By externalizing this knowledge, MCP allows it to persist and evolve independently of individual interactions.\\ \newline
Procedural memory captures action sequences or skills that can be applied across multiple situations. MCP supports procedural memory through prompt primitives that encode standard procedures or approaches, making them available for consistent application. This capability helps agents maintain consistent methodologies across interactions while adapting to specific circumstances.\\ \newline
Working memory maintains task-relevant information during ongoing activities. MCP enhances working memory through tool primitives that can track state information, maintain intermediate results, or monitor progress toward goals. This external augmentation of working memory helps agents manage complex tasks that exceed their internal context capacity.\\ \newline
By supporting these diverse memory types through standardized interfaces, MCP enables more human-like context awareness and memory utilization. Agents can remember relevant past interactions, apply accumulated knowledge, follow consistent procedures, and maintain awareness of ongoing tasks—capabilities that are essential for truly autonomous operation.\\
\subsubsection{Comparative Analysis with Previous Approaches}
Compared to previous approaches to context management, MCP offers several distinctive advantages:\\ \newline
Standardization vs. custom implementations: Earlier approaches typically implemented context management through custom, application-specific mechanisms with limited interoperability. MCP establishes standard interfaces and communication patterns that work across different applications, models, and environments. This standardization reduces implementation complexity, enables component reuse, and facilitates ecosystem development.\\ \newline
External vs. embedded context: Many previous approaches attempted to embed all relevant context within the model's input, leading to context window limitations and inefficient use of model capacity. MCP externalizes context management, storing information outside the model and retrieving it selectively based on relevance. This approach overcomes context window constraints while enabling more sophisticated context prioritization.\\ \newline
Structured vs. unstructured context: Earlier approaches often treated context as unstructured text, limiting the model's ability to distinguish between different types of contextual information. MCP provides structured primitives (prompts, resources, tools) with defined semantics, enabling more precise context management and utilization. This structure helps models understand the nature and purpose of different contextual elements.\\ \newline
Persistent vs. ephemeral context: Previous approaches frequently struggled with maintaining context across multiple interactions or sessions, leading to discontinuity in extended tasks. MCP provides mechanisms for persistent context storage and retrieval, enabling continuity across interaction boundaries. This persistence is essential for tasks requiring extended reasoning chains or collaborative work.\\ \newline
Interoperable vs. isolated tools: Earlier tool-using agent frameworks typically implemented custom, framework-specific approaches to tool integration. MCP standardizes tool definitions, invocation patterns, and result handling, enabling tools to be shared across different frameworks and environments. This interoperability expands the range of capabilities available to agents while reducing implementation overhead.\\ \newline
These comparative advantages position MCP as a significant advancement over previous context management approaches. By addressing the fundamental limitations that have constrained agent capabilities, MCP enables more effective autonomous operation and collaboration in multi-agent systems.\\ \newline
The Model Context Protocol represents a critical foundation for the multi-agent system architecture described in the following section. By providing standardized mechanisms for context management and sharing, MCP enables more effective coordination and collaboration among diverse agent types, addressing one of the key challenges identified in previous research.\\ \newline
